% 来源：Zhu et al., HyMEM, Findings ACL 2026；PDF 第1页（印刷页11287）
\papersection{HyMEM}{电脑操作智能体的混合自演化记忆}{lanshuheng9-oss}

% 来源：HyMEM，PDF 第1--3页，表1；下例为机制说明
\begin{frame}{网页操作中的经验表示与更新}
    \large\key{网页操作需要视觉细节，任务推进也会改变所需经验。}\par
    \vspace{5mm}\normalsize
    \begin{itemize}\setlength{\itemsep}{4mm}
        \item 输入：任务描述、当前截图、历史成功操作轨迹。
        \item 输出：指导智能体选择下一步网页操作。
        \item 例：搜索商品转入结账后，保留预算，更新操作指导。
    \end{itemize}
    \takeaway{混合表示解决“存什么”；自演化解决“何时更新”。}
    \note{约50秒。HyMEM面向网页操作智能体。任务输入包含描述、当前截图和历史成功轨迹，输出用于指导下一步动作。文本摘要容易遗漏界面细节；任务从搜索商品进入结账时，所需指导也会变化。预算等长期约束需要保留，当前操作经验则需要刷新。因此，方法同时研究经验怎样表示，以及怎样随任务更新。}
\end{frame}

% 来源：HyMEM，PDF 第3--5页，式(1)、图1、Structured Retrieval
\begin{frame}{混合记忆：策略与轨迹表示}
    \begin{columns}[T,onlytextwidth]
        \column{.52\textwidth}
        \centering
        \begin{tikzpicture}[node distance=5mm,box/.style={draw=accent,rounded corners,align=center,text width=53mm,minimum height=9mm,font=\small},>=Stealth]
            \node[box] (q) {任务文本＋当前截图};
            \node[box,below=of q] (r) {检索5条种子＋5条图邻居};
            \node[box,below=of r] (m) {离散指导指令＋连续轨迹嵌入};
            \node[box,below=of m] (a) {智能体执行下一步动作};
            \draw[->] (q)--(r);\draw[->] (r)--(m);\draw[->] (m)--(a);
        \end{tikzpicture}
        \column{.44\textwidth}
        \textbf{一条经验的三个层次}
        \[v_i=(c_i,A_i,m_i)\]
        \begin{itemize}\setlength{\itemsep}{3mm}
            \item $c_i$：可复用的高层策略。
            \item $A_i$：动作、界面及领域属性。
            \item $m_i$：8个连续轨迹嵌入。
        \end{itemize}
        \par\vspace{3mm}CLIP 用于检索；连续轨迹嵌入供智能体使用。\par
    \end{columns}
    \note{约85秒。系统根据任务和截图检索五条种子，再沿图扩展五条邻居，形成当前需要的经验。每条经验包含高层策略、动作与界面属性，以及八个连续轨迹嵌入。高层策略帮助规划，连续表示保留截图和动作中的细节。CLIP用于检索相似经验，轨迹编码器生成供智能体使用的连续表示。两种表示共同进入决策过程。}
\end{frame}

% 来源：HyMEM，PDF 第4--5页，Memory Update、Working Memory Refresh；图1
\begin{frame}{长期经验与工作记忆的更新}
    \begin{columns}[T,onlytextwidth]
        \column{.47\textwidth}
        \key{全局：跨任务积累经验}\par\vspace{3mm}
        \begin{itemize}\setlength{\itemsep}{3mm}
            \item ADD：出现新策略或属性。
            \item MERGE：同一策略的互补证据。
            \item REPLACE：同条件下更优的轨迹。
        \end{itemize}
        \par\vspace{3mm}先检索候选，再判断新轨迹提供了哪些信息。\par
        \column{.47\textwidth}
        \key{局部：跟随当前任务阶段}\par\vspace{3mm}
        \begin{enumerate}\setlength{\itemsep}{3mm}
            \item 比较动作前后观测，检测阶段变化。
            \item 保留长期目标及约束，删除过时指导。
            \item 重新检索，刷新文本与连续嵌入。
        \end{enumerate}
    \end{columns}
    \takeaway{全局改变长期经验库；局部更新当前执行上下文。}
    \note{约85秒。更新分为全局和局部。全局把新的成功轨迹与长期经验比较：新策略可以添加，互补信息可以合并，同条件下更优的轨迹可以替换旧轨迹。局部则关注当前任务阶段，在动作后比较观测，判断是否需要保留目标、移除过时指导并重新检索。两层更新都依赖模型判断，误合并和误刷新是需要检查的错误类型。}
\end{frame}

% 来源：HyMEM，PDF 第5--6页，\S4.1；第7页，表2表注
\begin{frame}{实验设置与编码器训练}
    \begin{tabularx}{\linewidth}{@{}p{25mm}Y@{}}
        \toprule
        项目 & 作者实验条件 \\
        \midrule
        三个基准 & WebVoyager、Multimodal-Mind2Web、MMInA \\
        记忆来源 & 2883条成功轨迹；主图1858个节点，超过100万条边 \\
        模型与接口 & 7B/8B视觉语言模型；ReAct、SOM及定位回退 \\
        训练 & Q-Former＋LoRA；约1.2\%参数可训练，通常1个epoch \\
        评价 & 任务成功率；Overall按可用领域取平均 \\
        \bottomrule
    \end{tabularx}
    \par\vspace{3mm}前两个基准评价截图与操作轨迹；MMInA评价最终回答。\par
    \note{约65秒。实验覆盖三个网页操作基准，记忆来自GUIAct、Mind2Web训练集及作者生成的成功轨迹。连续记忆编码器使用Q-Former和LoRA训练，约百分之一点二的参数可训练。评价采用任务成功率：前两个基准根据任务和截图轨迹判分，MMInA比较最终回答。记忆图、视觉输入、训练和动作接口共同构成这组实验条件。}
\end{frame}

% 来源：HyMEM，PDF 第7页，表2；图2为不同分析设置，不与本表混合
\begin{frame}{混合记忆的任务成功率}
    \textbf{任务成功率（\%），按可用领域取平均}\par\vspace{3mm}
    \begin{tabularx}{\linewidth}{@{}Y*{4}{>{\centering\arraybackslash}p{17mm}}@{}}
        \toprule
        基础模型 & 无记忆 & 离散 & 连续 & HyMEM \\
        \midrule
        Qwen2.5-VL-7B & 12.5 & 17.0 & 21.7 & \key{35.0} \\
        Qwen3-VL-8B & 21.9 & 21.3 & 23.3 & \key{33.8} \\
        UI-TARS-1.5-7B & 17.1 & 15.4 & 28.7 & \key{31.2} \\
        \bottomrule
    \end{tabularx}\par
    \vspace{4mm}
    \key{12.5\%到35.0\%：提高22.5个百分点。}\par\vspace{2mm}
    Qwen3 的离散记忆为21.3\%，低于无记忆时的21.9\%。

    \note{约90秒。Qwen2.5这一行中，无记忆时成功率为百分之十二点五，混合记忆达到百分之三十五，提高二十二点五个百分点。Qwen3的离散记忆略低于无记忆，连续记忆与混合记忆则更高。论文还分别分析全局演化和局部刷新：Amazon任务中，全局配置从四十八点八提高到六十一，局部配置提高到五十六点一。这是两组独立对照，可以分别观察两层维护的作用。}
\end{frame}

% 来源：HyMEM，PDF 第3--7页；右侧为本节阅读分析与待验证假设
\begin{frame}{方法总结与课设启发}
    \begin{columns}[T,onlytextwidth]
        \column{.47\textwidth}
        \textbf{贡献与局限}\par\vspace{3mm}
        \begin{itemize}\setlength{\itemsep}{3mm}
            \item 离散策略与视觉细节互补。
            \item 跨任务积累经验，任务内刷新上下文。
            \item 依赖训练、标签与模型判断，存在额外调用成本。
        \end{itemize}
        \column{.47\textwidth}
        \textbf{对事件记忆检索的启发}\par\vspace{3mm}
        \begin{itemize}\setlength{\itemsep}{3mm}
            \item 可研究种子检索与结构扩展的分工。
            \item 固定预算，比较覆盖、噪声及成本。
            \item 区分文本证据检索与多模态轨迹使用。
        \end{itemize}
    \end{columns}

    \note{约75秒。HyMEM把可复用策略与视觉轨迹结合，并把跨任务积累和任务内刷新分开。代价是额外训练、轨迹数据和模型调用。对事件记忆检索，可借鉴种子检索与结构扩展的分工，在固定预算下比较证据覆盖、噪声和成本。这样可以分析哪些信息来自直接命中，哪些来自结构关联。下面介绍关注事实使用的MetaMem。}
\end{frame}
